% Consequences of the effective residue estimate, /root/cm_effectivity.
% IDs NC01--NC04. Uses theorem/proposition/proof environments and amsmath/amssymb.
% Bibliography keys are supplied in cm/bibliography.tex.

\section{Effective CM class numbers and CM-point orbits}

Throughout this section, $D_E$ denotes the positive absolute discriminant of
$E$, and $\kappa_E=\operatorname*{Res}_{s=1}\zeta_E(s)$.
We use the quadratic Dirichlet GRH theorem supplied as Alp\"oge's theorem
and carried in ZETA\_TWO IA05.  No statement in this section uses Masullo's
claim, the general individual Colmez conjecture, or GRH for a nonquadratic
Hecke character.  The residue estimate used from NX02 is
\begin{equation}\label{eq:NC-residue-input}
 \kappa_E>\frac{1}{173[ E:\mathbf Q]!\log D_E}
 \qquad([E:\mathbf Q]\geq2).
\end{equation}
For normal $E$, NX01 permits the factorial to be replaced by $1$.

\begin{theorem}[NC01: explicit Siegel bounds for CM class numbers]
\label{NC01}
Let $E$ be any CM field of degree $2g$, let $F=E^+$, and put
\[
 c_g=\frac{2}{173(2\pi)^g(2g)!}
       \begin{cases}
       1,&g=1,\\
       \bigl(4(g-1)/e\bigr)^{g-1},&g\geq2.
       \end{cases}
\]
Then
\begin{equation}\label{eq:NC01-log}
 h_E\ \geq\ \frac{h_E}{h_F}\ >\
 c_g\frac{D_E^{1/4}}{(\log D_E)^g}.
\end{equation}
Consequently, for every explicitly given $\varepsilon>0$,
\begin{equation}\label{eq:NC01-epsilon}
 h_E>c_g\left(\frac{e\varepsilon}{g}\right)^g
                 D_E^{1/4-\varepsilon}.
\end{equation}
If $E$ is normal, the same statements hold with $c_g$ multiplied by
$(2g)!$.  In particular, for $H\geq1$, the condition $h_E\leq H$ implies
\begin{equation}\label{eq:NC01-cutoff}
 D_E<\left\{\frac{H}{c_g(e/(8g))^g}\right\}^{8}.
\end{equation}
For imaginary quadratic fields one retains the natural, stronger exponent:
\begin{equation}\label{eq:NC01-quadratic}
 h_E>\frac{\sqrt{D_E}}{173\pi\log D_E}
 \geq\frac{e\varepsilon}{173\pi}D_E^{1/2-\varepsilon}
 \qquad(\varepsilon>0).
\end{equation}
The constants are deliberately coarse and are completely explicit.
The stronger large-discriminant quadratic estimates in
ZETA\_TWO IA02--IA03 remain available; they are not reproved here.
\end{theorem}

\begin{proof}
The relative analytic class number formula, in the normalization of
\cite[\S3, equation (21), p.~1200]{NC-Louboutin2001}, is
\begin{equation}\label{eq:NC01-relative}
 \frac{h_E}{h_F}=
 \frac{Q_Ew_E}{(2\pi)^g}
 \sqrt{\frac{D_E}{D_F}}\frac{\kappa_E}{\kappa_F},
 \qquad Q_E\in\{1,2\},\quad w_E\geq2.
\end{equation}
This formula accounts for the regulator ratio
$R_E/R_F=2^{g-1}/Q_E$; replacing a bound for $h_ER_E$ by a bound for
$h_E$ without this step would be invalid.
For $g\geq2$, Louboutin's explicit upper bound
\cite[Theorem~1(1), p.~1195]{NC-Louboutin2001} gives
\[
 \kappa_F\leq
 \left(\frac{e\log D_F}{2(g-1)}\right)^{g-1}.
\]
Combining this with \eqref{eq:NC-residue-input} gives the stronger
field-dependent bound
\begin{equation}\label{eq:NC01-strong}
 \frac{h_E}{h_F}>
 \frac{2}{173(2\pi)^g(2g)!\log D_E}
 \sqrt{\frac{D_E}{D_F}}
 \left(\frac{2(g-1)}{e\log D_F}\right)^{g-1}.
\end{equation}
The discriminant identity
$D_E=D_F^2N_{F/\mathbf Q}(\mathfrak d_{E/F})$ implies
$D_F\leq D_E^{1/2}$.  Hence the square root in
\eqref{eq:NC01-strong} is at least $D_E^{1/4}$, while
$\log D_F\leq\tfrac12\log D_E$.  This proves
\eqref{eq:NC01-log} for $g\geq2$.  For $g=1$ use
$F=\mathbf Q$, $\kappa_F=1$, and
$D_E^{1/2}\geq D_E^{1/4}$ directly in
\eqref{eq:NC01-relative}.  Finally $h_F\geq1$.

For $t>0$, the maximum of $t^g e^{-\varepsilon t}$ occurs at
$t=g/\varepsilon$ and equals $(g/(e\varepsilon))^g$.
With $t=\log D_E$, this proves \eqref{eq:NC01-epsilon}.
Taking $\varepsilon=1/8$ proves \eqref{eq:NC01-cutoff}.
Normality changes only the residue constant, as stated.
For \eqref{eq:NC01-quadratic}, a quadratic field is normal and has
$R_E=1$, so NX01 and the ordinary class number formula give
$h_E=w_E\sqrt{D_E}\kappa_E/(2\pi)>
\sqrt{D_E}/(173\pi\log D_E)$.
The same calculus inequality with $g=1$ gives its second inequality.
\end{proof}

\paragraph{Arithmetic-geometric meaning of NC01.}
Here $h_E=|\operatorname{Pic}(\operatorname{Spec}\mathcal O_E)|$.
The unit logarithm determinant in $K_1(\mathcal O_E)$ is the regulator;
\eqref{eq:NC01-relative} cancels this determinant against its counterpart
for $\operatorname{Spec}\mathcal O_F$.  On the generic fibres,
$\zeta_E/\zeta_F$ is the $L$-function of the genuine Artin summand
$\operatorname{Ind}_{G_F}^{G_{\mathbf Q}}\chi_{E/F}$ in
$H^0(\operatorname{Spec}E)$, where $\chi_{E/F}$ is the quadratic
character of $G_F$.  The passage through the normal closure in NX02
controls a possible exceptional quadratic constituent, not every zero
of this Artin $L$-function.  Thus NC01 changes the effectivity of the
class-number estimate without asserting the full RH of $\zeta_E$.

\begin{theorem}[NC02: an explicit logarithmic CM height bound]
\label{NC02}
For $g\geq1$, define
\[
 a_g=2^{g-1}\left(\frac{75}{2}+16g!\right),
 \qquad b_g=4g(2^g-1),\qquad H_g(D)=a_g\log D+b_g.
\]
Let $A/\overline{\mathbf Q}$ have dimension $g$ and complex
multiplication by the full ring $\mathcal O_E$ of a CM field of degree
$2g$.  Let $h_F$ be the stable height in Faltings' original normalization.
Then
\begin{equation}\label{eq:NC02-height}
 h_F(A)\leq H_g(D_E).
\end{equation}
This assertion holds for every CM type, including imprimitive types.
\end{theorem}

\begin{proof}
We spell out the constants in the proved argument of
\cite[Proposition~5.3 and its proof, pp.~19--20]{NC-Wei2021}.
Use first the height $h_W$ defined with the squared archimedean norm
$(2\pi)^{-g}\int|\omega\wedge\overline\omega|$.
This is the normalization of Yuan--Zhang
\cite[\S1.1, p.~534; Theorem~1.1, p.~535]{NC-YuanZhang2018}
and Wei.  Gaudron--R\'emond
\cite[\S2.3, pp.~351--352]{NC-GRPeriods2014} give
\begin{equation}\label{eq:NC02-normalization}
 h_W(A)=h_F(A)+\frac g2\log\pi.
\end{equation}
Their Theorem~1.1, p.~345, has positive left hand side and upper bound
$16h_F(A)+49g$, so it implies $h_F(A)>-49g/16$.
In particular $h_W(A)>-4g$.

Stark's Lemma~9, pp.~145--146, applied with $[F:\mathbf Q]=g$, says
that a real zero of $L(s,\chi_{E/F})=\zeta_E(s)/\zeta_F(s)$ in
\[
 1-\frac{1}{16g!\log D_E}\leq s<1
\]
would descend to an imaginary quadratic subfield of $E$.
Alp\"oge's quadratic Dirichlet GRH excludes this.  The only possible
zero in the wider square of side $1/(4\log D_E)$ is real and simple
(Stark Lemma~3; also \cite[\S1, p.~3]{NC-Wei2021}), so
$C_{\mathrm{zero}}(g)=16g!$ satisfies the hypothesis of Wei's
Proposition~5.3.

Let $\Phi_0$ be the type of $A$.  The averaged Colmez identity is
\[
 2^{-g}\sum_\Phi h_W(E,\Phi)=
 -\frac12\frac{L'(0,\chi_{E/F})}{L(0,\chi_{E/F})}
 -\frac14\log(D_E/D_F).
\]
Wei's estimate in the proof of Proposition~5.3 bounds the right hand
side above by
\[
 \frac12\left\{
 \left(\frac{75}{2}+16g!\right)\log D_E
 +\frac g2\bigl(\psi(1/2)-\log\pi\bigr)\right\},
\]
where $\psi$ is the logarithmic derivative of the gamma function.
Isolate the summand $\Phi_0$ and use $h_W(E,\Phi)>-4g$ for the
other $2^g-1$ summands.  Since
$\psi(1/2)-\log\pi=-\gamma-\log(4\pi)<0$, discarding the last
term gives $h_W(A)\leq H_g(D_E)$.
Equation \eqref{eq:NC02-normalization} then gives
\eqref{eq:NC02-height}.  No individual Colmez formula was used.
\end{proof}

\paragraph{Geometric meaning of NC02.}
The height is the normalized arithmetic degree of the metrized Hodge
line $\det(e^*\Omega^1_{\mathcal A/\mathcal O_K})$ on an abelian
scheme after a field extension giving good reduction.
Averaging the CM types expresses that degree in terms of the logarithmic
derivative of a weight-zero Artin $L$-function.  The polarizable
weight-one CM realization $H^1(A)$ is present through this Hodge line;
its own infinite-order Hecke $L$-function is not assumed to satisfy GRH.
NC02 strengthens the averaged-height consequence recorded after
ZETA\_TWO IA05 and does not settle the individual Colmez identity
which IA05 explicitly leaves as a separate input.

\subsection{Explicit constants for the orbit statements}

For each integer $r\geq1$, let $c_r,a_r,b_r$ have the values just
specified and define
\begin{align*}
 \alpha_r&=2^{10}r^3,& B_r&=(14r)^{64r^2},&
 \eta_r&=\frac{1}{16\alpha_r},\\
 U_r&=a_r+\frac{b_r+1}{\log3},&
 q_r&=2+\frac{r}{2\alpha_r},&
 s_r&=2\cdot3^{4r^2},\\
 f_r&=\frac{c_r^{1/(2\alpha_r)}}
              {2^{r/\alpha_r}B_rU_r^2},&&&\\
 \ell_r&=\min\left\{1,\,
           f_r\left(\frac{e\eta_r}{q_r}\right)^{q_r}\right\}.&&&
\end{align*}
Every constant here is positive and given by elementary operations.
The isogeny estimate to be used is
\cite[Theorem~1.4, p.~2060; definition of $\kappa(A)$, p.~2058]
{NC-GRPol2014}: two isogenous abelian varieties $A,A'$ of dimension $r$
defined over a number field of degree $k$ have an isogeny in each
direction of degree at most
\begin{equation}\label{eq:NC-GR}
 \bigl(B_rk\max\{h_F(A),\log k,1\}^{2}\bigr)^{\alpha_r}.
\end{equation}
The theorem permits the isogeny over any extension on which the two
varieties are isogenous.  In the applications below there is already
an isogeny over the stated field.

\begin{theorem}[NC03: primitive full-ring CM orbits]
\label{NC03}
Let $A/\overline{\mathbf Q}$ be geometrically simple of dimension $g$
with $\operatorname{End}(A)=\mathcal O_E$ for a CM field $E$ of degree
$2g$.  If $A$ and all its endomorphisms are defined over $K$, then
\begin{equation}\label{eq:NC03-def}
 [K:\mathbf Q]\geq\ell_g D_E^{\eta_g}.
\end{equation}
For a polarized $A$ with field of moduli $\mathbf Q(A)$,
\begin{equation}\label{eq:NC03-moduli}
 [\mathbf Q(A):\mathbf Q]\geq
 \frac{\ell_g}{s_g}D_E^{1/(16384g^3)}.
\end{equation}
This is an effective version of Tsimerman's Theorems~1.2 and~4.2.
\end{theorem}

\begin{proof}
Put $D=D_E$, $k=[K:\mathbf Q]$, and fix the primitive CM type $\Phi$
of $A$.  The usual CM ideal construction gives $h_E$ isomorphism
classes in $S(E,\Phi)$; the quotient by the kernel of an integral ideal
$I\subset\mathcal O_E$ has isogeny degree $N(I)$, and every class
arises in this way.  The precise facts used are recorded in
\cite[\S2.2, p.~382, proof of Proposition~2.2]{NC-Tsimerman2018}.
Because all endomorphisms of $A$ are defined over $K$, their common
ideal kernels and these quotients are defined over $K$ as well.
For each quotient $A_I$, apply \eqref{eq:NC-GR} to $A,A_I$ over $K$.
It gives an isogeny $A\to A_I$ of degree at most
\[
 N=\bigl(B_gk\max\{H_g(D),\log k,1\}^{2}\bigr)^{\alpha_g}.
\]
All homomorphisms in this isogeny class commute with the fixed
$E$-action: after fixing one $E$-linear isogeny, the rational Hom space
is a rank-one module over the commutative algebra
$\operatorname{End}^0(A)=E$.  Thus the isogeny just obtained is among
those classified by integral ideals.  There are consequently at least
$h_E$ integral ideals of norm at most $N$.

For any number field of degree $2g$ and any $N\geq1$,
\[
 \#\{I\subset\mathcal O_E:N(I)\leq N\}
 \leq N^2\zeta_E(2)
 \leq N^2\zeta(2)^{2g}<2^{2g}N^2.
\]
The middle inequality follows by bounding each rational-prime Euler
factor by $(1-p^{-2})^{-2g}$.  Hence $h_E<2^{2g}N^2$.
If $k>D$, \eqref{eq:NC03-def} follows from $\ell_g\leq1$ and
$\eta_g<1$.  Assume therefore $k\leq D$.  As $D\geq3$,
\[
 \max\{H_g(D),\log k,1\}\leq U_g\log D.
\]
The lower bound NC01 and the last ideal count give
\[
 c_g\frac{D^{1/4}}{(\log D)^g}
 <2^{2g}(B_gkU_g^2(\log D)^2)^{2\alpha_g}.
\]
Taking $2\alpha_g$-th roots and solving for $k$ yields
\[
 k>f_g\frac{D^{1/(8\alpha_g)}}{(\log D)^{q_g}}.
\]
For $t>0$, the maximum of $t^{q_g}e^{-\eta_gt}$ is
$(q_g/(e\eta_g))^{q_g}$.  Using $t=\log D$ proves
\eqref{eq:NC03-def} by the definition of $\ell_g$.
Finally Tsimerman's Lemma~4.1, p.~384, gives a field of definition of
$A$ and all its endomorphisms of degree at most
$s_g[\mathbf Q(A):\mathbf Q]$.  (The lemma is formulated with
unpolarized field of moduli and its extension defines every
polarization.  Taking the compositum with the polarized field of
moduli preserves the bound $s_g$ on the relative degree.)
Apply \eqref{eq:NC03-def} to that field.
\end{proof}

\paragraph{Which quantitative steps changed.}
The analytic inputs to the orbit argument are now the explicit
class-number lower bound NC01 and the logarithmic height bound NC02.
The ideal count is elementary.  The isogeny estimate
\eqref{eq:NC-GR} has specified constants, so it leaves no ineffective
constant in NC03.  Its exponent is intentionally sacrificed to keep
one uniform, directly evaluable expression.  This is an orbit estimate;
no numerical degree bound for arbitrary special subvarieties is being
claimed from it.

\begin{theorem}[NC04: all CM points, including nonmaximal orders]
\label{NC04}
For $g\geq1$ define the following further explicit constants:
\begin{align*}
 \beta_g&=16\alpha_g,\\
 M_g&=\max_{1\leq r\leq g}\ell_r^{-16\alpha_r},\\
 A_g^*&=\max_{1\leq r\leq g}a_r,
 &C_g^*&=\max_{1\leq r\leq g}b_r,\\
 V_g&=1+g\{A_g^*(\log M_g+\beta_g)+C_g^*\},\\
 T_g&=M_g^g(B_gV_g^2)^{4g\alpha_g},
 &P_g&=28g\alpha_g=28672g^4.
\end{align*}
Let $x\in\mathcal A_g(\overline{\mathbf Q})$ be a CM point, let
$A_x$ be its principally polarized abelian variety, and let
$R_x=Z(\operatorname{End}_{\overline{\mathbf Q}}A_x)$, an order in a
product of CM fields.  Then
\begin{equation}\label{eq:NC04-orbit}
 |\operatorname{Gal}(\overline{\mathbf Q}/\mathbf Q)\cdot x|
 \geq s_g^{-1}T_g^{-1/P_g}\,|\operatorname{disc}R_x|^{1/P_g}.
\end{equation}
In particular the full Galois-orbit assertion of
\cite[Theorem~5.2, p.~385]{NC-Tsimerman2018} has an explicit exponent
and an explicit multiplicative constant under Alp\"oge's theorem.
\end{theorem}

\begin{proof}
Put $d=[\mathbf Q(x):\mathbf Q]$ and $A=A_x$.
By the same Lemma~4.1 and the compositum just described,
choose $K$ over $\mathbf Q(x)$ with
$k=[K:\mathbf Q]\leq s_gd$ such that $A$ and every geometric
endomorphism are defined over $K$.

R\'emond's maximal-order replacement theorem
\cite[Theorem~1.1, pp.~1--2 of the author preprint]{NC-Remond2017}
provides a $K$-isogenous abelian variety $B/K$ whose endomorphism
ring is a maximal order.  This theorem works over the same base
field: no extension depending on the index of the original order
is needed.  All geometric rational endomorphisms of $B$ are defined
over $K$, by transport through the $K$-isogeny.  The integral ones
are therefore defined over $K$ too.

Repeated application of R\'emond's Theorem~1.2, p.~2, decomposes
$B$ over $K$ as a product of simple factors $B_j$ with maximal
endomorphism rings.  Every geometric endomorphism of each factor is defined over $K$,
by its inclusion into and projection from $B$.  Such a factor is
geometrically simple: applying the same Theorem~1.2 over
$\overline{\mathbf Q}$ to a geometric proper abelian subvariety
would give a nontrivial idempotent endomorphism.  This idempotent
is defined over $K$, contradicting simplicity over $K$.
As $A$ is CM, each factor has dimension $r_j\leq g$,
$\operatorname{End}(B_j)=\mathcal O_{E_j}$ for a CM field of degree
$2r_j$, and primitive CM type (\cite[\S7, paragraph preceding Lemma~7.5,
pp.~17--18 of the preprint]{NC-Remond2017}).  There are at most $g$ factors,
counting repetitions, with $\sum_jr_j=g$.

NC03 in its field-of-definition form implies, for every $j$,
\[
 D_{E_j}\leq(k/\ell_{r_j})^{16\alpha_{r_j}}
          \leq M_g k^{\beta_g}.
\]
The stable Faltings height is additive on products
(\cite[\S2.3, p.~352]{NC-GRPeriods2014}).  NC02 thus gives
\begin{align*}
 h_F(B)&\leq\sum_j\{a_{r_j}\log D_{E_j}+b_{r_j}\}\\
 &\leq g\{A_g^*\log M_g+C_g^*\}
          +gA_g^*\beta_g\log k.
\end{align*}
It follows that
\[
 \max\{h_F(B),\log k,1\}
 \leq V_g(1+\log k)\leq V_gk,
\]
where the last inequality holds for $k\geq1$.
The isogeny estimate \eqref{eq:NC-GR}, applied with $B$ as the
height input, supplies an isogeny $\varphi:B\to A$ of degree
\begin{equation}\label{eq:NC04-degree}
 N\leq(B_gV_g^2)^{\alpha_g}k^{3\alpha_g}.
\end{equation}

Identify the centres of the rational endomorphism algebras by
$\varphi$.  Let $\mathcal O$ be the integral closure of $\mathbf Z$
in this centre.  Since $\operatorname{End}(B)$ is maximal,
its centre is $\mathcal O$: adjoining central algebraic integers
would otherwise give a strictly larger order.  Also
$R_x\subset\mathcal O$, because an integral endomorphism acts on
an integral homology lattice and is integral over $\mathbf Z$.
The kernel of $\varphi$ is killed by its order $N$, so there is an
isogeny $\psi:A\to B$ with $\psi\varphi=[N]_B$.
For $u\in\mathcal O$ the endomorphism
$\varphi u\psi$ of $A$ is central and represents $Nu$ under our
identification.  Thus
\[
 N\mathcal O\subset R_x\subset\mathcal O,
 \qquad [\mathcal O:R_x]\leq N^{\operatorname{rank}_{\mathbf Z}\mathcal O}
                         \leq N^{2g}.
\]
The discriminant of the product order $\mathcal O$ is the product
of the discriminants of its field factors, one for each simple
isogeny component, retaining separate components even if their
fields are isomorphic.  There are at most $g$ of these.  The discriminant-index formula
and \eqref{eq:NC04-degree} give
\begin{align*}
 |\operatorname{disc}R_x|
 &=|\operatorname{disc}\mathcal O|[\mathcal O:R_x]^2\\
 &\leq M_g^gk^{g\beta_g}N^{4g}\\
 &\leq M_g^g(B_gV_g^2)^{4g\alpha_g}
             k^{g\beta_g+12g\alpha_g}
   =T_gk^{P_g}\leq T_g(s_gd)^{P_g}.
\end{align*}
Solving for $d$, which is precisely the Galois-orbit cardinality,
proves \eqref{eq:NC04-orbit}.
\end{proof}

\paragraph{Geometric meaning and exact scope of NC03--NC04.}
The points lie in the moduli scheme $\mathcal A_g/\mathbf Q$.
Isogeny preserves the rational CM motive and all its $\ell$-adic
realizations, while replacing the endomorphism order changes an integral
lattice.  NC04 quantifies that lattice change by
$N\mathcal O\subset R_x$ and the discriminant-index identity.
The Hodge line supplies the height used to bound the isogeny.
Hence these results make the CM-orbit input to Tsimerman's proof
effective, including the centre of a nonmaximal endomorphism ring.
They use the abelian category and its endomorphisms, not a presumed
motivic splitting for arbitrary varieties.  The weight-one non-CM
factors in ZETA\_TWO ZT06 remain a different question.

\paragraph{The distinction from a full effective Andr\'e--Oort algorithm.}
NC04 supplies all constants in the arithmetic orbit lower bound requested
here.  The complete proof in \cite[\S6.3, Theorem~6.1, and \S6.5,
p.~388]{NC-Tsimerman2018} also uses a subpolynomial counting theorem and
an o-minimal reduction from special subvarieties to special points.
Those cited steps, as stated there, supply no algorithm for their
constants or finite families; such an algorithm does not follow from
NC04 alone.  This identifies the scope of the present proof, not a
claim that effective replacements cannot exist.  Indeed, for powers
of the modular curve there already is such a replacement:
Binyamini \cite[\S1.2, Theorem~1, preprint p.~2; \S1.4, pp.~3--4]
{NC-Binyamini2020} gives effective bounds for maximal special
subvarieties once the possible exceptional imaginary quadratic field
is removed by Dirichlet GRH.  That specialization is effective under
the supplied Alp\"oge theorem.  The two nonarithmetic effectivity
issues for the original o-minimal proof are listed separately in the
same \S1.4.

\paragraph{Archive and programme locators.}
These quantitative class-number and moduli-height estimates are not
statements found in EGA, SGA, FGA, or the supplied Grothendieck scans.
Their direct programme targets are ZETA\_TWO IA05 and its intake
\S6(1)(a), and the appended CM class-number discussion in
INTAKE\_RESEARCH\_SA06\_20261120, \S4, rows B.4a--B.4e, restated
in ZETA\_TWO IAQ-008--012.  The supplied intake does not identify a
cote-and-scan for B.4a--B.4e; none is invented here.  The explicit
archive connection specified in ZETA\_TWO's intake \S5.4 is
cote~15, scan~70 for SA06 A.8a, concerning quadratic
class-number/supersingular questions.  The related row A.8c has
neither a separate cote-and-scan locator nor a definition of its
ring $R$ in the supplied intake.  The sharper quadratic estimates
remain ZETA\_TWO IA02--IA03.
The cohomological determinant interpretation is the one already fixed
in ZETA\_TWO ZT08: SGA~6, Exp.~V, \S2 and (7.1.1), and the programme's
cote~60, scans~3--9.  These locators concern the framework, not a
claim that the new numerical inequalities occur in the archive.
